\subsection{Web Application Testing}

The web application testing plan verifies both the user-facing interface and the backend API that the web app, mobile app, and microcontroller share. 
Tests run locally with mocked authentication, database responses, and network calls so behavior can be checked without live hardware or a production database. 
Frontend tests use a component test library to render pages and simulate user actions. The backend tests invoke API handlers directly and assert on status codes, response bodies, and error handling. 
All automated tests live under the \texttt{test/} directory in the web-server source tree and are executed with a single test runner command.

\subsubsection{Frontend and User Interface}

Frontend tests focus on the two panels: the Account Settings page and the Setup Guide page. These tests confirm that the application shell renders both sections, 
that unauthenticated users see a sign-in form, and that authenticated users see account information loaded from the API.
For the Settings page, tests cover the full signed-out and signed-in flows. When no session exists, the page shows a loading state and then the sign-in form.
Invalid credentials display an error message from the authentication service. After a successful sign-in, the page lists the user's groups and devices, 
handles API load failures gracefully, and updates when the session changes. Users can sign out, which clears displayed account data. 
Security actions are also verified: account deletion requires confirmation and only proceeds when the user accepts, and cancellation leaves the session intact.
The Setup Guide page is tested at the layout level: both the Account and Setup Guide headings appear together so users can read setup instructions alongside account management. 
Because the guide is largely static content, testing emphasizes correct rendering rather than dynamic data flows.

\subsubsection{Backend API and Authentication}

The web server also hosts REST endpoints used by the mobile app and the microcontroller. A separate set of tests validates server-side logic in isolation. 
Each test calls an API handler with a mock request and inspects the HTTP response, using stand-in database and auth clients so tests do not depend on external services.
Authentication tests verify that user requests without a valid Bearer token are rejected, that device ingestion requests require valid device credentials, 
and that invalid or replayed tokens are not accepted. These checks protect endpoints that read or write user and device data.
Group-management tests cover listing, creating, joining, renaming, and deleting groups, as well as updating calibration values and managing members (including role changes and removal). 
Tests include both successful paths and failure cases such as invalid invite codes, duplicate membership, missing permissions, and database errors.
Device tests verify listing a user's registered devices and the BLE registration flow used during mobile provisioning, including linking a new device to an existing group and handling missing or invalid input.
The sensor ingestion endpoint is tested for correct handling of POST requests, device authentication, weight-to-level conversion using calibration data, storage of readings, 
and optional low-water push notifications. Profile endpoints are tested for reading and updating user information, signing out all active sessions, and deleting an account.
Supporting utility modules, such as push-notification delivery and shared group logic, are also tested so edge cases (empty inputs, network failures, malformed responses) do not propagate into production handlers.
